\chapter{Receiver Architectures}
\label{ch:receivers}

A radio receiver's job is to select one wanted signal from among the
enormous number of other signals and noise present at its antenna,
convert it to a usable form, and produce audio (or digital data) with
adequate signal-to-noise ratio and freedom from interference. This
chapter surveys the major receiver architectures used from the earliest
days of radio to the present.

\section{Key Receiver Performance Concepts}

\begin{itemize}
\item \textbf{Sensitivity}: the weakest signal a receiver can usefully
detect, limited ultimately by internally generated noise (thermal noise
in resistive elements and semiconductor devices), quantified by the
receiver's \textbf{noise figure} (how much noisier the receiver is than
an ideal noiseless one, in dB).
\item \textbf{Selectivity}: the receiver's ability to reject signals on
frequencies close to the wanted one, set largely by its filter
bandwidth and shape (Chapter~\ref{ch:rlc-resonance}).
\item \textbf{Dynamic range}: the span, in dB, between the weakest
signal the receiver can detect and the strongest signal it can handle
without unacceptable distortion (\textbf{overload} or
\textbf{intermodulation distortion}, where two or more strong nearby
signals mix together inside a nonlinear stage to create spurious
``phantom'' signals).
\item \textbf{Image rejection}: the ability to reject an unwanted
frequency that a mixing stage would otherwise convert to the same output
frequency as the wanted signal (Section~\ref{sec:superhet} below).
\end{itemize}

\section{Tuned Radio Frequency (TRF) Receivers}

\bookfig[0.55]{receptor_trf.png}{Block diagram of a tuned radio
frequency (TRF) receiver.}

The earliest practical receiver architecture amplifies the incoming
signal directly at its original (radio) frequency through one or more
stages of RF amplification, each individually tuned to the wanted
frequency, before demodulating it. TRF receivers are conceptually
simple but have a fundamental weakness: keeping several separate tuned
stages aligned to exactly the same frequency as the operator tunes
across the band is mechanically difficult, and selectivity (which
depends on $Q$, Chapter~\ref{ch:rlc-resonance}) is hard to make
uniformly good across a wide tuning range, since a fixed-$Q$ circuit's
absolute bandwidth changes with frequency.

\section{Direct Conversion Receivers}

\bookfig[0.55]{receptor_convdirecta.png}{Block diagram of a direct
conversion receiver.}

A \textbf{direct conversion} receiver mixes the incoming RF signal
directly with a local oscillator (LO) set to (nearly) the same
frequency as the wanted signal, producing an output directly at audio
frequency (the mixing process is developed fully in
Chapter~\ref{ch:modulation}). This is simpler than a superheterodyne
design (fewer stages, no IF filter needed) and is popular in simple and
homebrew equipment, but suffers from being susceptible to
\textbf{audio-frequency image response} (signals above and below the LO
frequency both fold into the same audio output, requiring phasing
techniques to separate sidebands) and generally offers less overall
selectivity than a well-designed superheterodyne.

\section{The Superheterodyne Receiver}
\label{sec:superhet}

\bookfig[0.6]{receptor_superhet.png}{Block diagram of a superheterodyne
receiver.}

The \textbf{superheterodyne} (``superhet''), developed by Edwin
Armstrong and still the dominant architecture (now often implemented
digitally rather than with discrete analog stages), solves the TRF
alignment problem by converting every incoming signal to a fixed
\textbf{intermediate frequency (IF)}, regardless of the tuned frequency,
using a mixer and a tunable local oscillator:
\[
\boxed{f_{\text{IF}} = |f_{\text{signal}} - f_{\text{LO}}|}
\]
Because the IF is fixed, the bulk of the receiver's gain and
selectivity can be built once, at a single frequency, using
well-optimized fixed-tuned filters (often crystal or ceramic filters)
--- giving consistently good, easily reproducible selectivity across the
entire tuning range, the key advantage over a TRF design.

\subsection{The Image Frequency Problem}

Because the mixer output depends only on the \emph{difference} between
signal and LO frequency, \emph{two} input frequencies produce the same
IF output: the wanted signal at $f_{\text{LO}} + f_{\text{IF}}$ (or
$f_{\text{LO}} - f_{\text{IF}}$) and an unwanted \textbf{image
frequency} on the opposite side of the LO, offset by exactly
$2f_{\text{IF}}$ from the wanted signal:
\[
f_{\text{image}} = f_{\text{signal}} \pm 2f_{\text{IF}}.
\]

\begin{examplebox}[Finding an image frequency]
A receiver tuned to 14.200\,MHz uses a 9.000\,MHz IF and high-side
injection ($f_{\text{LO}} = f_{\text{signal}} + f_{\text{IF}} =
23.200\,\text{MHz}$). Find the image frequency.

\emph{Solution.}
\[
f_{\text{image}} = f_{\text{LO}} + f_{\text{IF}} = 23.200 + 9.000 = 32.200\,\text{MHz}.
\]
(Equivalently, $f_{\text{signal}} + 2f_{\text{IF}} = 14.200 + 18.000 =
32.200\,\text{MHz}$.) A signal at 32.200\,MHz, if not rejected before the
mixer, would appear at the same 9.000\,MHz IF as the wanted 14.200\,MHz
signal.
\end{examplebox}

\begin{keyidea}[Fighting the image]
Image rejection is provided by front-end filtering \emph{before} the
mixer. A \emph{higher} IF pushes the image frequency further away from
the wanted signal (easier to filter with a simple front-end filter), but
makes it harder to build a narrow, high-selectivity IF filter at that
higher frequency (Chapter~\ref{ch:rlc-resonance}: absolute bandwidth for
a given $Q$ grows with frequency). This trade-off is why many
modern receivers use a high first IF for good image rejection, followed
by conversion to a lower second IF for good adjacent-channel
selectivity --- a \textbf{dual (or multiple) conversion} superheterodyne.
\end{keyidea}

\section{The Superregenerative Receiver}

\bookfig[0.5]{receptor_superreg.png}{A simplified superregenerative
detector.}

A \textbf{superregenerative} receiver deliberately drives a single RF
amplifier stage into and out of oscillation at a rapid (typically
ultrasonic) rate, exploiting the huge effective gain available right at
the threshold of oscillation to achieve remarkable sensitivity from a
very small number of components. It trades selectivity and
linearity for extreme simplicity and low cost, and remains common in
inexpensive, short-range devices (some remote controls and simple
telemetry receivers) despite being essentially obsolete for
general-purpose amateur communication receivers.

\section{Software-Defined Radio (SDR)}

A \textbf{software-defined radio} performs some or all of the
traditional analog stages --- filtering, mixing, demodulation --- in
digital signal processing after an analog-to-digital converter (ADC)
samples the incoming RF (or an early IF) directly. Modern amateur SDR
receivers commonly use a \textbf{direct sampling} or
\textbf{quadrature-sampling-detector (QSD)} front end feeding a
general-purpose or dedicated digital signal processor, implementing all
subsequent filtering, mixing, AGC, and demodulation as software
algorithms rather than discrete analog circuit stages. This gives
extraordinary flexibility (mode, bandwidth, and filter shape can be
changed instantly in software) and, in modern designs, performance that
matches or exceeds the best conventional analog superheterodynes, at the
cost of requiring fast, high-dynamic-range ADCs and significant digital
processing power. SDR concepts underlie most currently manufactured
amateur transceivers, whether marketed as ``SDR'' or built with a more
conventional-looking analog front panel.

\section{Automatic Gain Control (AGC)}

Because received signal strength can vary over many orders of magnitude
(weak DX signals to strong local stations), nearly all practical
receivers include \textbf{automatic gain control}: a feedback loop that
senses the receiver's output level and automatically reduces gain for
strong signals (and restores it for weak ones), keeping the audio output
level roughly constant and protecting later stages from overload. AGC
response speed (fast vs.\ slow ``attack'' and ``release'' times) is
usually adjustable and matters for how comfortably different modes ---
CW, SSB voice, or the sudden noise bursts common on HF --- sound to the
operator.

\begin{quickreview}
\item Receiver quality is described by sensitivity (noise figure),
selectivity, dynamic range, and image rejection.
\item TRF receivers amplify directly at RF; hard to align across a wide
tuning range.
\item Direct conversion mixes straight to audio; simple, but prone to
image response unless phasing techniques are used.
\item Superheterodyne receivers convert every signal to a fixed IF for
consistent selectivity; the trade-off is the image frequency,
$f_{\text{image}} = f_{\text{signal}} \pm 2f_{\text{IF}}$.
\item SDR performs filtering/demodulation digitally after early
analog-to-digital conversion, offering great flexibility and, in modern
designs, top-tier performance.
\item AGC automatically adjusts receiver gain to keep output level
manageable across widely varying signal strengths.
\end{quickreview}

\begin{practiceproblems}
\item A superhet uses a 455\,kHz IF to receive 7.150\,MHz with high-side
injection. Find the LO frequency and the image frequency.
\item Explain why a higher IF improves image rejection but can worsen
adjacent-channel selectivity for a fixed filter $Q$.
\item Explain, conceptually, why a direct conversion receiver needs a
phasing technique to reject one sideband, while a superheterodyne
typically does not for the same purpose.
\item A receiver's noise floor is $-130\,\text{dBm}$ and the strongest
signal it can handle without unacceptable distortion is $-10\,\text{dBm}$.
Find its dynamic range in dB.
\item Explain, in your own words, one practical advantage SDR receivers
have over fixed analog-filter superheterodynes.
\end{practiceproblems}
